\SourceLeaf{404}%
\par
Entwickelt man diese beiden Kosinus, dividiert dann durch \(\cos\tfrac12\varphi\cos\tfrac12\varphi'\) und bezeichnet \(\operatorname{tang}\tfrac12\varphi\) .mit \(t\), \(\operatorname{tang}\tfrac12\varphi'\) mit \(t'\)
\[
\text{und}\quad\frac{a\sin A}{r-a\cos A}=\alpha,\quad\frac{r+a\cos A}{r-a\cos A}=\beta,\quad\text{so wird}
\]
\[
\text{I.}\qquad t=\tfrac{1-\alpha t'}{\alpha-\beta t'}
\]
\par
Ganz ähnlich erhält man, wenn man
\[
\operatorname{tang}\tfrac12\varphi''=t''\quad\frac{a'\sin A'}{r'-a'\cos A'}=\alpha',\quad\frac{r+a'\cos A'}{r-a'\cos A'}=\beta'
\]
setzt,
\[
\text{II.}\qquad t'=\tfrac{1-\alpha't''}{\alpha'-\beta't''}\ \text{usw.}
\]
\par
Daran sieht man, dass sich ebenso viele Gleichungen ergeben, wie das Vieleck Seiten hat. Verbindet man sie miteinander, erhält man schließlich eine quadratische Gleichung für \(t\).
\par\medskip\Needspace{5\baselineskip}
\pdfbookmark[2]{VI.}{carnot-zusatz-vi}
\centerline{\textbf{VI.}}\nobreak\smallskip
\textsc{Aufgabe}. Drei Kreise sind nach Lage und Größe gegeben. Man soll einen vierten Kreis zeichnen, der alle drei berührt.

\textit{Lösung}. Man lege die senkrecht aufeinander stehenden Achsen durch den Mittelpunkt des einen Kreises und bezeichne die Abstände von diesen Geraden wie folgt
\[
\begin{array}{rl}
\text{des Mittelpunkts des zweiten Kreises}\quad .\ .\ .\ .\ .&a,\quad b\\
\text{des dritten Kreises}\quad .\ .\ .\ .\ .&a',\quad b'\\
\text{des gesuchten}\quad .\ .\ .\ .\ .&x,\quad y
\end{array}
\]
\par
Der Abstand des Mittelpunkts des ersten Kreises vom Mittelpunkt des gesuchten Kreises sei \(=z\). Dann ist
\begin{list}{}{\setlength{\leftmargin}{3em}\setlength{\labelwidth}{2.5em}\setlength{\labelsep}{.5em}}
\item[\(z-c\)] der Abstand des Mittelpunkts des zweiten Kreises vom Mittelpunkt des gesuchten Kreises,
\item[\(z-c'\)] der Abstand des Mittelpunkts des dritten Kreises vom Mittelpunkt des gesuchten Kreises,
\end{list}
